% Stage 8 section framing — 2.9 Action of Radiation Operators on Angular Spectra
\section{Action of Radiation Operators on Angular Spectra}
\label{sec:ch29}
Reciprocity in Section~\ref{sec:ch28} couples ports and probes. Section~\ref{sec:ch29} asks how
radiation operators act on angular spectra---the language of pattern decomposition, near-field
scanning, and plane-wave synthesis \cite{jackson1999}. Every qualified
far-zone pattern is first a weighted sum of directional contributions: each term reports how much
export reaches a unit direction $\widehat{\mathbf{k}}$ on the sphere after outgoing selection
\cite{stratton1941}. That decomposition precedes any lateral Fourier transform;
plane-wave spectra are one coordinate system on the same directional budget, not a separate
physical layer \cite{jackson1999}. Maxwell propagation from compact support cannot populate
every direction with equal strength: outgoing selection removes incoming admixture; evanescent
sheets carry reactive content that decays with distance and does not survive the near-to-far map;
finite apertures integrate over unresolved directions. The accessible amplitude
$a(\widehat{\mathbf{k}})$ is therefore a compressed image of impressed structure, not a complete
coordinate on $\mathcal{J}_\Omega$ \cite{stratton1941}. On the bench, a spherical
near-field scanner records tangential $\E$ and $\H$ on a closed or planar surface at height $h$
above the device under test; the angular-spectrum chart synthesizes outgoing plane-wave trials, back-propagates to
the scan plane, and accepts residuals only after scan pitch, bandlimit, and probe correction are
stated \cite{hansen1988}. The sequence scan--synthesize--bandlimit--report is what the
theory encodes once radiation has been selected \cite{stratton1941}.




Section~\ref{sec:ch28} closed the reciprocity chain
Eq.~\eqref{eq:ch2.section28-reciprocity-chain}: Lorentz reciprocity (Subsection~\ref{sec:ch281}),
modal orthogonality and coupling (Subsection~\ref{sec:ch282}), and radiation--reception
interchange (Subsection~\ref{sec:ch283}) \cite{stratton1941}. That program fixed exchange
symmetry on \(\operatorname{ran}\mathcal{R}_\omega\) and on port coordinates after projection.
Section~\ref{sec:ch29} develops filter action in plane-wave coordinates: how
\(\mathcal{R}_\omega\), \(\Pi_{\mathrm{rad}}\), and observation maps retain or suppress
continuous angular content on \(S^{2}\), and how near-zone fields map to far-zone amplitudes
\(a(\widehat{\mathbf{k}})\) \cite{stratton1941}. Preview expansion
Eq.~\eqref{eq:ch2.angular-spectrum} from Subsection~\ref{sec:ch224} named the lateral
spectral template; here that leading template is developed to locked statements compatible with outgoing
criterion Eq.~\eqref{eq:ch2.outgoing-plane-wave-criterion} and far amplitude
Eq.~\eqref{eq:ch2.far-field-amplitude-integral} \cite{stratton1941}. Explicit basis change to
vector harmonics is developed in \ref{ch:03} Subsection~3.8.1; harmonic normalization belongs to
Subsection~3.5.2 with citation of Subsection~\ref{sec:ch282} \cite{harrington2001}.
Spectral decomposition in Subsection~\ref{sec:ch261} and continuous representation
Eq.~\eqref{eq:ch2.exterior-continuous-spectral-integral} fixed modal and plane-wave coordinates
on \(\mathcal{F}_{\mathrm{rad}}\); null-space consequences in Subsection~\ref{sec:ch274} and
angular budget Eq.~\eqref{eq:ch2.angular-spectrum-composition-law} classified which directions
\(a(\widehat{\mathbf{k}})\) may carry export \cite{stratton1941}. Reciprocity from
Section~\ref{sec:ch28} constrains bilinear coupling on observable coordinates but does not
replace angular filters: a reciprocal antenna may still have suppressed sheets
Eq.~\eqref{eq:ch2.suppressed-angular-sheet-def} and finite observation rank
Eq.~\eqref{eq:ch2.angular-spectrum-observation-rank} \cite{stratton1941}. Treating
\(\mathcal{S}_\omega[\Jimp]\) near sources as if it were already \(a(\widehat{\mathbf{k}})\)
repeats the regime error class of Section~\ref{sec:ch01.2.4} \cite{stratton1941}.
Section~\ref{sec:ch29} separates plane-wave angular spectra, retained and suppressed
components, near-to-far transformation, and far-field angular $\mathcal{E}_{\mathrm{acc}}^{(N)}$ without re-opening
Green construction from Section~\ref{sec:ch23} \cite{stratton1941}.
Let \(\omega>0\) be fixed, \(\Omega_s=\operatorname{supp}\Jimp\) be compact, and outgoing pairs
lie in \(\mathcal{F}_{\mathrm{SF}}\) Eq.~\eqref{eq:ch2.admissible-field-space}. Section
~\ref{sec:ch29} organizes analysis in four stages:
\begin{equation}
\label{eq:ch2.section29-angular-spectrum-chain}
\text{plane-wave spectrum}
\ \xrightarrow{\ \mathcal{R}_\omega,\Pi_{\mathrm{rad}}\ }
\text{retained / suppressed}
\ \xrightarrow{\ \text{near-to-far}\ }
a(\widehat{\mathbf{k}})
\ \xrightarrow{\ \text{probe rank}\ }
\mathcal{K}_{\mathrm{acc}},
\end{equation}
where the first arrow is Subsection~\ref{sec:ch291}, the second is Subsection~\ref{sec:ch292},
the third is Subsection~\ref{sec:ch293}, and the fourth is Subsection~\ref{sec:ch294}
\cite{stratton1941}. Equation~\eqref{eq:ch2.section29-angular-spectrum-chain}
expresses Eq.~\eqref{eq:ch2.section27-nullspace-chain} and
Eq.~\eqref{eq:ch2.section28-reciprocity-chain} at the angular level: null-space and reciprocity
constraints from Sections~\ref{sec:ch27}--\ref{sec:ch28} bound which \(a(\widehat{\mathbf{k}})\)
are physically meaningful; Section~\ref{sec:ch29} develops the operators that enforce those
bounds on plane-wave coordinates \cite{stratton1941}. Transport criteria from
Section~\ref{sec:ch01.2.3} and angular observables from Section~\ref{sec:ch01.2.1} enter
Subsection~\ref{sec:ch294} by reference, not by re-derivation \cite{stratton1941}.
Modal representation, not separate physics (Sections~2.6 and~2.9). Spectral and angular decompositions organize information about \(\mathcal{R}_\omega[\Jimp]\);
they do not define angular structure \cite{stratton1941}. Observable angular structure remains
primary: modes, plane-wave sheets, and harmonic labels are coordinates on export, not substitutes
for Maxwell flux or qualified transport \cite{maxwell1873}. A nonzero modal coefficient or
angular amplitude reports surviving source-generated content after propagation and outgoing
selection; renaming that content does not create new physical degrees of freedom
\cite{harrington2001}.
The primary objects are radiation map
Eq.~\eqref{eq:ch2.nullspace-operator-targets}, outgoing plane-wave criterion
Eq.~\eqref{eq:ch2.outgoing-plane-wave-criterion}, and continuous angular amplitude
\(a(\widehat{\mathbf{k}})\) from Eq.~\eqref{eq:ch2.exterior-continuous-spectral-integral}
\cite{stratton1941}. Secondary targets include propagating projectors
Eqs.~\eqref{eq:ch2.spectral-projectors}--\eqref{eq:ch2.rad-null-on-evanescent},
observation composition Eq.~\eqref{eq:ch2.reconstruction-observation-map}, accessible sector
Eq.~\eqref{eq:ch2.angular-accessibility-sector}, and composition law
Eq.~\eqref{eq:ch2.angular-spectrum-composition-law} \cite{stratton1941}. Section~\ref{sec:ch29}
develops filters on the same outgoing branch fixed by
Eqs.~\eqref{eq:ch2.outgoing-radiation-refinement}--\eqref{eq:ch2.outgoing-resolvent-unique};
incoming sheets remain excluded from \(\mathcal{F}_{\mathrm{rad}}\)
\cite{harrington2001}.
Plane-wave angular spectrum and radiation operators (leading template). Subsection~\ref{sec:ch291} supplies the definitive account of plane-wave angular spectra
and the action of \(\mathcal{R}_\omega\circ\Pi_{\mathrm{prop}}\) on spectral coefficients (Eq.~\eqref{eq:ch2.plane-wave-radiation-operator} in Subsection~\ref{sec:ch291}): \begin{equation}
\widetilde{\E}_{\mathrm{rad}}(k_x,k_y)
=
\mathcal{F}_{\mathrm{rad}}\big[\Pi_{\mathrm{rad}}\,\mathcal{R}_\omega[\Jimp]\big](k_x,k_y),
\qquad
k_z=k_z(k_t)\ \text{from Eq.~\eqref{eq:ch2.outgoing-plane-wave-criterion}},
\end{equation}
with \(\mathcal{F}_{\mathrm{rad}}\) the lateral Fourier transform on outgoing admissible fields
\cite{stratton1941}. Equation~\eqref{eq:ch2.plane-wave-radiation-operator}
develops Eq.~\eqref{eq:ch2.angular-spectrum} by inserting radiation selection;
evanescent sheets satisfy \(\widetilde{\E}_{\mathrm{rad}}=0\) on export through
Eq.~\eqref{eq:ch2.rad-null-on-evanescent} \cite{stratton1941}. Rigged-space normalization of
plane-wave modes belongs to \ref{ch:03} Subsection~3.1.3 \cite{harrington2001}.
Subsection~\ref{sec:ch292} partitions angular content into retained, suppressed, and
observation-filtered sheets:
\begin{equation}
\label{eq:ch2.retained-suppressed-filter}
a(\widehat{\mathbf{k}})
=
a_{\mathrm{ret}}(\widehat{\mathbf{k}})\,\chi_{\mathcal{K}_{\mathrm{acc}}}(\widehat{\mathbf{k}})
\;+\;
a_{\mathrm{sup}}(\widehat{\mathbf{k}})\,\chi_{\mathcal{K}_{\mathrm{sup}}}(\widehat{\mathbf{k}}),
\qquad
a_{\mathrm{filt}}=\mathcal{O}_{\mathrm{meas}}[a],
\end{equation}
using accessible sector Eq.~\eqref{eq:ch2.angular-accessibility-sector} and suppressed sheet
Eq.~\eqref{eq:ch2.suppressed-angular-sheet-def} from Subsection~\ref{sec:ch274}
\cite{stratton1941}. Equation~\eqref{eq:ch2.retained-suppressed-filter}
is the Section~\ref{sec:ch29} target for Subsection~\ref{sec:ch292} partition; geometric
constraints on radiative spectra (Subsection~\ref{sec:ch264}) and observation rank
Eq.~\eqref{eq:ch2.angular-spectrum-observation-rank} compose independently with
Eq.~\eqref{eq:ch2.angular-spectrum-composition-law} \cite{stratton1941}. Null modes annihilate
\(a(\widehat{\mathbf{k}})\) through Eq.~\eqref{eq:ch2.angular-spectrum-nullspace-law} without
altering near reactive content \cite{harrington2001}.
Near-to-far transformation (leading template). Subsection~\ref{sec:ch293} records the operator map from near-zone admissible fields to far
angular amplitudes:
\begin{equation}
\label{eq:ch2.near-to-far-transform}
a(\widehat{\mathbf{k}})
=
\lim_{r\to\infty}
\mathcal{H}_{\mathrm{ff}}\big[\mathcal{R}_\omega[\Jimp]\big](r,\widehat{\mathbf{k}}),
\end{equation}
with \(\mathcal{H}_{\mathrm{ff}}\) the far-zone extractor tied to
Eq.~\eqref{eq:ch2.far-field-amplitude-integral} and the far-field approximation of source integrals (Subsection~\ref{sec:ch252}) \cite{stratton1941}.
Equation~\eqref{eq:ch2.near-to-far-transform} does not promote evanescent storage to
export: \(\Pi_{\mathrm{ev}}\mathcal{S}_\omega[\Jimp]\) is annihilated before
Eq.~\eqref{eq:ch2.near-to-far-transform} when \(\mathcal{R}_\omega\) is unchanged
\cite{stratton1941}. regime qualification Eq.~\eqref{eq:ch2.regime-component-criterion} qualifies when the
limit defines transport data \cite{harrington2001}.
Subsection~\ref{sec:ch294} develops probe rank on \(S^{2}\) importing Section~\ref{sec:ch01.2.1}
angular observables and Section~\ref{sec:ch01.2.3} transport criteria by reference
\cite{stratton1941}:
\begin{equation}
\label{eq:ch2.far-angular-accessibility-preview}
\widehat{\mathbf{k}}\in\mathcal{K}_{\mathrm{acc}}
\ \Longleftrightarrow\
a(\widehat{\mathbf{k}})\ \text{observable after}\ \mathcal{M}_\omega
\ \text{and qualified by}\ \mathcal{D}_{J,\mathrm{reg}}=1,
\end{equation}
upgrading angular $\mathcal{E}_{\mathrm{acc}}^{(N)}$ on the far shell (Subsection~\ref{sec:ch264}) and composing with
Eq.~\eqref{eq:ch2.angular-spectrum-composition-law} \cite{stratton1941}. Equation
\eqref{eq:ch2.far-angular-accessibility-preview} is the Section~\ref{sec:ch29} completion target;
explicit plane-wave filters and measurement qualification is developed in Subsection~\ref{sec:ch294}
\cite{harrington2001}. Reciprocity in radiation and reception (Subsection~\ref{sec:ch283}) constrains
coupling on \(\mathcal{K}_{\mathrm{acc}}\) but does not enlarge it \cite{stratton1941}.
\begin{equation}
\label{eq:ch2.angular-action-composition}
\Jimp\in\ker\mathcal{R}_\omega
\ \Longrightarrow\
a(\widehat{\mathbf{k}})=0
\ \text{on export},
\qquad
\mathfrak{C}_\omega\ \text{symmetric on}\ \mathcal{K}_{\mathrm{acc}},
\end{equation}
importing Eq.~\eqref{eq:ch2.angular-spectrum-nullspace-law} and reciprocity in radiation and reception (Subsection~\ref{sec:ch283}) \cite{stratton1941}. Modal
coordinates Eq.~\eqref{eq:ch2.radiation-map-modal-sum} and plane-wave coordinates
Eq.~\eqref{eq:ch2.exterior-continuous-spectral-integral} describe the same
\(\operatorname{ran}\mathcal{R}_\omega\) on exterior domains; Section~\ref{sec:ch29} fixes
plane-wave-side filters, while \ref{ch:03} Subsection~3.8.1 supplies basis
change law \cite{stratton1941}. Finite-rank observation
Eq.~\eqref{eq:ch2.reconstruction-rank-cap} may remove angular content even when
\(a(\widehat{\mathbf{k}})\neq 0\) on \(\mathcal{F}_{\mathrm{rad}}\) \cite{stratton1941}.
Several constructions deliberately remain outside Section~\ref{sec:ch29}. Vector spherical
harmonic completeness and normalization are developed in \ref{ch:03}; Section~\ref{sec:ch29}
records plane-wave radiation action without tabulating \(Y_\ell^m\)
\cite{stratton1941}. Non-radiating catalogs remain in Section~\ref{sec:ch27};
reciprocity constraints remain in Section~\ref{sec:ch28} \cite{stratton1941}. Section~\ref{sec:ch29}
uses Fredholm structure, geometric constraints, and null-space structure from
Sections~\ref{sec:ch26}--\ref{sec:ch27} \cite{harrington2001}.
Regime logic from Section~\ref{sec:ch01.2.4} is unchanged: near-zone reactive dominance does not
create far-zone angular transport without \(\mathcal{D}_{J,\mathrm{reg}}=1\) and outgoing
qualification \cite{stratton1941}. Angular-transport claims must compose with
Eq.~\eqref{eq:ch2.section29-angular-spectrum-chain}, angular budget
Eq.~\eqref{eq:ch2.angular-spectrum-composition-law}, and regime qualification
Eq.~\eqref{eq:ch2.regime-operator-screen} \cite{harrington2001}.
